\section{Baseline Experiments}\label{sec:baseline}
\subsection{Objective and common evaluation conditions}
Notebook 03 establishes whether the selected predictors support improvement over simple load-history rules before hyperparameter search. All h24 candidates use the same \num{7105} training hours ending on 31 October and the same 720 November validation hours. The five default ML models use the 17-column set D defined in \cref{sec:features}; the baselines access their required historical columns on the same eligible timestamps. This comparison measures performance in November, rather than accuracy averaged across seasonal regimes.

Naive-24h and Naive-168h copy the preceding day's and preceding week's corresponding hour. Their usefulness follows from the daily and weekly dependence in \cref{fig:acf}. MeanProfile instead pools the training observations associated with each hour and weekday. Writing $I_{a,d}=\{i\in\mathcal T:a_i=a,\ d_i=d\}$ for the training subset, its prediction is
\begin{equation}
\widehat y_t^{\mathrm{profile}}=\begin{cases}
|I_{a_t,d_t}|^{-1}\displaystyle\sum_{i\in I_{a_t,d_t}}y_i,&|I_{a_t,d_t}|>0,\\
|\mathcal T|^{-1}\displaystyle\sum_{i\in\mathcal T}y_i,&\text{otherwise}.
\end{cases}\label{eq:meanprofile}
\end{equation}
The profile and its fallback are fitted only to training observations. The seasonal rules have no fitted coefficients, although their histories update as measured loads become available.

The additional RF-P1 configuration has 30 trees and seven candidate features per split, with unrestricted depth and the split/leaf minima specified in \cref{sec:method}. Its role is to test a literature-inspired parameter configuration within this project's target and validation design. It does not establish replication of an external paper's scores.

\subsection{November results}
\Cref{tab:baselines,fig:baseline_comparison} show the saved comparison. XGB has the lowest validation RMSE, \num{2161.4} kW, and also the lowest MAE and MAPE, \num{1598.4} kW and \num{2.42}\%, respectively. Against Naive-24h at \num{3076.4} kW RMSE, its RMSE reduction is \num{29.74}\%. RF-P1, SVR and RF follow at \num{2192.0}, \num{2195.2} and \num{2227.5} kW. These relatively close November results do not establish statistically significant differences.
\begin{table}[tbp]\centering\small
\caption{November validation, h24: seasonal baselines, default regressors and RF-P1 (T04). RMSE/MAE are in kW; MAPE is percent.}\label{tab:baselines}
\setlength{\tabcolsep}{4pt}
\begin{tabular}{lrrrrr}
\toprule
Model & Train RMSE & Val. RMSE & Val. MAE & MAPE & R$^2$ \\
\midrule
XGB & \num{812.4} & \num{2161.4} & \num{1598.4} & \num{2.42} & \num{0.9760} \\
RF-P1 & \num{805.0} & \num{2192.0} & \num{1620.3} & \num{2.50} & \num{0.9753} \\
SVR & \num{1567.7} & \num{2195.2} & \num{1601.1} & \num{2.52} & \num{0.9753} \\
RF & \num{794.4} & \num{2227.5} & \num{1616.7} & \num{2.47} & \num{0.9745} \\
LR & \num{2670.5} & \num{2288.7} & \num{1734.0} & \num{2.76} & \num{0.9731} \\
DT & \num{0.0} & \num{2992.4} & \num{2190.0} & \num{3.34} & \num{0.9540} \\
Naive-24h & \num{3497.5} & \num{3076.4} & \num{2080.7} & \num{3.27} & \num{0.9514} \\
Naive-168h & \num{4686.4} & \num{3202.9} & \num{2425.8} & \num{3.79} & \num{0.9473} \\
MeanProfile & \num{9628.3} & \num{9369.6} & \num{8560.8} & \num{13.50} & \num{0.5493} \\
\bottomrule\end{tabular}
\end{table}

\begin{figure}[tbp]\centering
\includegraphics[width=\linewidth]{F11_level1_vs_baselines.png}
\caption{November h24 RMSE and MAPE for seasonal baselines, default regressors and RF-P1 (F11). The vertical references mark Naive-24h. Rankings and baseline improvements apply to this validation month.}\label{fig:baseline_comparison}
\end{figure}

The default decision tree illustrates why a statement that a model simply ``beats the baseline'' is incomplete. Its validation RMSE of \num{2992.4} kW improves on Naive-24h by \num{2.73}\%, but its MAE, \num{2190.0} kW, exceeds the baseline's \num{2080.7} kW. Its MAPE is also worse: \num{3.34}\% against \num{3.27}\%. Thus fewer or smaller large errors can coexist with worse average absolute error. This metric-specific interpretation follows the current CSV rather than earlier prose that labels DT uniformly unsuccessful.

Naive-168h reaches \num{3202.9} kW RMSE, slightly worse than Naive-24h. Weekly alignment alone therefore does not compensate for changes in the load level in this month. MeanProfile reaches \num{9369.6} kW RMSE. A training-wide calendar average cannot adapt its absolute level as directly as a recent measured load. This is consistent with the observed seasonal variation, although the comparison does not isolate a single cause of profile error.

\subsection{Training fit and representative validation behavior}
DT fits its training observations to numerical precision, yet its validation error remains substantial. RF and XGB have training RMSEs of \num{794.4} and \num{812.4} kW, well below their validation errors. These gaps indicate that training fit alone would give a misleading ranking. Conversely, LR has validation RMSE \num{2288.7} kW against training RMSE \num{2670.5} kW. A lower validation error is possible when the validation period is easier; it does not imply that model fitting used validation labels.

\Cref{fig:baseline_week} displays 13--19 November, the second Monday-to-Sunday week selected by the notebook. The previous-day baseline copies Sunday's lower daytime demand into Monday and carries Saturday's pattern into Sunday. XGB corrects much of these mismatches while retaining the broad daily trajectory. Calendar and history features provide a plausible explanation, but this plot alone cannot assign the improvement to a particular predictor. The civil holiday on 18 November offers only a single illustrative date, not an independent estimate of a general holiday effect.
\begin{figure}[tbp]\centering
\includegraphics[width=\linewidth]{F11b_week_actual_vs_pred.png}
\caption{Actual load, default XGB forecasts and Naive-24h during 13--19 November (F11b). XGB was selected for this illustration by November validation RMSE. The week is part of model development and is not the December test.}\label{fig:baseline_week}
\end{figure}

The initial comparison supports proceeding to temporal validation: several regressors improve substantially on the seasonal baseline in November. It does not yet support selecting XGB for deployment, because a single autumn month omits the summer regime and provides no estimate of seasonal stability.
\evidence{Notebook 03 cells 6--19; \code{src/models.py}; \code{T04_level0_level1_wide.csv}; F11 and F11b. Numerical values refer to November validation unless explicitly identified as training metrics.}
